\begin{theorem}[Incident-core reduction]
Let $G=(V,E)$ be a finite loopless multigraph and let $H=G^{\circ}$ be its
incident core.  Then:
\begin{enumerate}
\item coercing the endpoints of every individual edge identity of $H$ back
to $V$ gives its endpoint pair in $G$;
\item every vertex of $H$ is incident with an edge;
\item $H$ is critical if and only if $G$ is critical;
\item $\chi'(H)=\chi'(G)$ and $\Delta(H)=\Delta(G)$;
\item $t(H)=t(G)$; and
\item if $H$ is elementary, then $G$ is elementary.
\end{enumerate}
The edge type of $H$ is exactly $E$, so all assertions preserve individual
parallel-edge identities.
\end{theorem}
\begin{proof}
The endpoint assertion is \noderef{IncidentCoreEndpoints}, and absence of
isolated core vertices is \noderef{IncidentCoreNoIsolated}.  Criticality
equivalence is \noderef{IncidentCoreCritical}.  The two numerical equalities
are \noderef{IncidentCoreChromaticIndex} and
\noderef{IncidentCoreMaximumDegree}; Tashkinov-order equality is
\noderef{IncidentCoreTashkinovOrder}.  Finally,
\noderef{IncidentCoreElementaryTransfer} transfers elementariness from $H$
back to $G$.  These are logically distinct properties proved from the same
incident-core construction, and their conjunction gives every clause of the
stated reduction.  Since \noderef{IncidentCore} uses the original edge type
$E$ throughout, the conjunction also retains each parallel identity exactly.
\end{proof}
